\section{A concrete path to complete Lean formalization}
\label{sec:formalization}

\subsection{Three levels that must remain distinct}

The report contains complete human arguments for the stated finite and
represented-output theorems. Their iPRA accounting identifies fixed PR
programs, bounded invariants, finite search limits, and the required
instances of unbounded input hypotheses. The root existence constructor
reuses the explicitly identified rational-descent lemma from the FTA
packet; an unchanged copy of its article is included in the archive.

The delivered Lean files are at two different implementation levels.
\path{BackwardDemand.lean} and \path{FiniteHausdorff.lean} prove
native Lean lemmas. \path{AggregateTargets.lean},
\path{RootSetTargets.lean}, and \path{NormTargets.lean} define
concrete algorithms and named propositions for their missing proofs.
A declaration of the form \texttt{def TheoremTarget : Prop := ...}
neither proves its proposition nor introduces it as an axiom.
All five files compile. The archive contains no admitted proof of a
main theorem and does not claim that these definitions constitute one.

The final logical-strength endpoint is different again:
\[
                  \texttt{Prf PRAAx [] sentence}.
\]
This type denotes a derivation in the repository's embedded arithmetic.
An arbitrary native Lean induction on \texttt{Prop} is not an instance
of the embedded theory's quantifier-free induction. In particular,
the native scheduling theorem deliberately accepts arbitrary predicates;
its iPRA application requires the separate boundedness and PR hypotheses
proved in Section~\ref{sec:backward-demands}.

\subsection{Definitions and theorem boundaries in the delivered modules}

\begin{longtable}{@{}p{0.24\textwidth}p{0.39\textwidth}p{0.30\textwidth}@{}}
\toprule
Module & Concrete content & Proof status\\
\midrule
\endhead
\texttt{BackwardDemand} & Backward schedule; forward evaluation;
parameter-changing recursion; finite query envelopes. & Schedule,
compilation equation, transfer, and finite-query theorems proved in
native Lean.\\[4pt]
\texttt{FiniteHausdorff} & Directed covers; Boolean checker; rational
metric; two-sided closeness; regular aggregate names. & Reflection,
composition, image, union, and triangle lemmas proved in native Lean.\\[4pt]
\texttt{AggregateTargets} & Rational vectors, exact minima, least-neighbor
chain, sampled membership, represented evaluation, and minimizers. &
Chain, distance, value, and approximate-witness goals are unproved
propositions.\\[4pt]
\texttt{RootSetTargets} & Gaussian rationals, polynomial operations,
descent clock, deflation, budgets, sampled root support and adequacy. &
Six targets: descent-to-root, product stability, factorization,
coherence, and two adequacy goals. All unproved.\\[4pt]
\texttt{NormTargets} & Rational polynomial lists, companion matrices, monic
reduction, determinant, norm, weighted degrees, points and subspaces. &
General and superelliptic norm, multiplicity, zero-count, and subspace
goals are unproved. Finite examples execute.\\
\bottomrule
\end{longtable}

\texttt{AggregateTargets.MemberZero} tests the rational samples of the
distance name at precision $k+1$ against $2^{-k}$. The least-neighbor
proof supplies the stronger $\tfrac58\,2^{-k}$ bound. Its optimization
targets specify both the value name and the complete two-parameter
point-name constructor. Thus an exact minimizer is not silently stored
as an assumed field of a structure.

\texttt{RootSetTargets} defines the supplied descent map's specification
in squared-norm arithmetic and implements its $J$-step iterator. The
remaining descent assumption has an exact upstream proof location.
Its adequacy targets use rational samples and bounded squared-distance
covers, avoiding comparisons of exact complex roots. Its integer
precision constants can be more generous than those in the paper:
choosing $b=4C$ instead of a binary length still guarantees $2^b\ge4C$.
This costs time but changes no theorem.

\texttt{NormTargets.PlaneData} includes all coefficients of a monic
equation in $Y$. Companion-matrix powers implement reduction even for
an arbitrary input $Y$-coefficient list. The main general target retains
the essential nonzero-norm hypothesis. The human coprime-degree theorem
proves nonvanishing automatically; the specialized Lean target states
that conclusion but does not yet prove it. These compiled definitions
use rational coefficients and rational points. Generalization to the
coded-field interface of the written theorem is a further formalization
step. The subspace target uses
explicit row-span witnesses and a nonzero annihilator certificate;
Gaussian elimination is the producer described in the human proof.

\subsection{The arithmetic compiler work, in dependency order}

The intended integration uses the pinned repository modules
\path{Syntax.lean}, \path{Deduction.lean}, and \path{PR.lean}
\cite{repo-core}. They already provide scoped terms and formulas,
substitution, intuitionistic deductions, PR function syntax, and
\texttt{PRA.qfInd}. The following sequence avoids a gap between
native correctness and the object-theory claim.
\begin{enumerate}
\item \textbf{Concrete arithmetic.} Introduce PR syntax for signed
integers, positive-denominator rationals, list lookup and folds,
polynomial convolution, bounded maxima and least search. Prove their
defining-equation specifications in \texttt{PRAAx}. Every malformed code
gets a default; correctness predicates assert the valid-input guards.
\item \textbf{Boolean compilation.} Implement the fixed-syntax
translation of Lemma~\ref{lem:booleanization}. Its universal and
existential fold invariants were given explicitly, so this step needs
no induction on an unbounded satisfaction predicate. The compiler can
run in Lean's metatheory; each output must be a genuine derivation.
\item \textbf{Schedule derivations.} Reify the two recursions of
Section~\ref{sec:backward-demands}. Derive the truncated-subtraction
identities and schedule equations from \texttt{PRA.precZero},
\texttt{PRA.precSucc}, equality substitution, and quantifier-free
induction. Build the syntax of the guarded invariant and apply
\texttt{PRA.qfInd}. The native proof gives the exact algebraic template.
\item \textbf{Finite metric algebra.} Prove rational triangle and
finite-minimum specifications by denominator clearing and prefix
induction. Translate \texttt{coverCheck\_correct} and the finite cover
combinators. Formalize the least-neighbor invariant and the explicitly
bounded geometric sum. Then implement the value and witness budgets.
\item \textbf{Rational polynomial analysis.} Prove squared-norm
multiplication, the power-difference identity, coefficient convolution,
synthetic division, and the rational descent lemma. The last is the
substantial inherited obligation. Compose its bounded iterator with
the fixed-length deflation invariant; then substitute $E(t)$ in the
finite product-stability proof and use the adequacy sample indices.
\item \textbf{Finite field algebra.} Prove monic reduction, evaluation
rows, determinant multiplication and row multilinearity, the
Vandermonde calculation, and synthetic removal of linear factors.
For the specialized theorem, formalize the unique-weight argument and
the unique cyclic-shift determinant term. The finite proof has no
irreducibility or divisor-theory dependency to import.
\item \textbf{Output adequacy and transport.} Universally introduce
the numerical parameters only after the finite derivations are built.
For named inputs, retain their validity hypotheses and instantiate them
at the explicitly calculated indices. The existing transport APIs between ordinary iPRA and
the ConciseIPRA scalar calculus can then be applied to an actual arithmetic
derivation; they do not produce that derivation from a native theorem.
\end{enumerate}

For example, in backward transfer the invariant formula is the syntax
of $t\le N\to c(t,h(t;N,K))=1$. The base is a derivation of its
substitution at zero. The step is a derivation in the context containing
that formula and the old context shifted by one variable. Instantiate
the local contract at $h(t+1)$, use the schedule equality, and introduce
the guard implication. After \texttt{PRA.qfInd}, eliminate the universal
quantifier at the term $N$ and rewrite $h(N)=K$. This specifies the
binders, substitutions, and final inference, not merely the shape of
the desired conclusion.

\clearpage
\subsection{Induction inventory}

\begin{longtable}{@{}p{0.23\textwidth}p{0.45\textwidth}p{0.25\textwidth}@{}}
\toprule
Construction & Finite invariant & Terminal resource\\
\midrule
\endhead
Backward transfer & $t\le N\to c(t,h(t))=1$ & $N$ forward steps and
$N$ backward steps.\\[3pt]
Least-neighbor chain & Current point lies in its finite list; accumulated
rational distance is at most the finite geometric sum. & $s+k+3$ as
largest observed list index for membership at $k$.\\[3pt]
Approximate factorization & Monicity and degree; coefficient bounds;
the exact partial telescoping identity; accumulated error. & $n$
divisions, each with its specified descent or grid bound.\\[3pt]
Norm construction & Each matrix column is the specified remainder;
partial determinant sum has its computed coefficient list. & $m^m$
candidate permutation words suffice.\\[3pt]
Norm zero count & Groups partition the input list; removed linear
factors multiply the current quotient to the original norm. & Input
list length and polynomial degrees.\\[3pt]
Subspace fibers & Pivot equations and selected free coordinate;
the annihilator is nonzero and annihilates each supplied row. & Matrix
row and column counts.\\
\bottomrule
\end{longtable}

The native axiom reports are included verbatim in the validation output.
The scheduling results use only \texttt{propext} and \texttt{Quot.sound}.
Pure relational cover composition uses no axioms. The native rational
triangle proof inherits \texttt{Classical.choice} from \texttt{Std}'s
rational-order lemmas. This explains why compiling the native rational
layer is not by itself a certificate of intuitionistic arithmetic
strength. The object development uses coded rational comparisons and
the explicit finite proofs instead.
